\scanpage{54}
% archive foliation: 53
\diagram{\begin{tikzcd}&X_s\arrow[r]&Y_s\\k(s)=k&T_s\end{tikzcd}}
\diagram{\begin{tikzcd}&X'_s\arrow[r,dash]&Y'_s\\k'&T'_s\arrow[u]\arrow[ur]\end{tikzcd}}
\margin{Nb de comp. connexes absolues de $T_s$: $\nu_s$
\[
\left\langle\begin{array}{l}
\nu=\sup\nu_s\\
\left\{\begin{array}{l}n=\deg T/S\\d=\dim.\ \text{fibres de }g\end{array}\right.\\
\lambda=\deg X/Y
\end{array}\right.
\]}
\deleted{\diagram{\begin{tikzcd}&\unsure{X_s|Y_s}\\\unsure{k}&T_s\arrow[u]\end{tikzcd}}}

\uline{Corollaire 1 bis} Sous les conditions du corollaire 1,
\deleted{Supposons $T/S$ fini,} \deleted{Alors} et supposons que \deleted{$\unsure{f}$} \inserted{le morphisme étale} $g:X\to Y$ soit propre de degré \deleted{$\unsure{d}$} $\lambda$.
Alors $f_*(g)_s:f_*(X/T)_s\to f_*(Y/T)_s$ est propre et de degré $\lambda^{\inserted{\nu_s}\deleted{\unsure{n}}}$ ($\nu_s=\deg_{\mathrm{sép}}T_s/k(s)$).

On reprend le raisonnement du cor. 1,
[\unsure{bloquer} cor. \unsure{3|bis} et cor. 1 en un seul énoncé]
Il suffit de prouver qu'il y a \uline{exactement} $\lambda^{\nu_s}$ points dans l'image inverse d'un pt de $f_*(Y/T)$.
\margin{Sous ces conditions, un morphisme \unsure{fini ét.|étale} de degré $\lambda$ est transformé en \unsure{un $=$ de gr.|un de gr.} $\lambda^\nu$}

\uline{Corollaire 3} (\deleted{$T/S$ fini}) \inserted{(Sous les cond. du cor. 1)} \deleted{\unsure{on déduit par le principe}} Conditions équivalentes:

(i) $\nu(s)$ est \inserted{loc.} constant

(ii) $X\mapsto f_*(X/T)$ transforme morphismes \deleted{finis} étales \uline{finis} en morphismes étales \uline{finis}

(iii) Si $T'$ est étale fini sur $T$, alors $f_*(T'/T)$ est étale fini sur \deleted{$S$}

(iv) $f_*(T\amalg T)\inserted{=\uline{\operatorname{Hom}}_S(T,S\amalg S)}$ \unsure{est} fini sur $S$ ?
\diagram{\begin{tikzcd}
(i)\arrow[d,Rightarrow]\\
(ii)\arrow[d,Rightarrow]\\
(iii)\arrow[d]\\
(iv)\arrow[uuu,bend left=55]
\end{tikzcd}}
